% !TeX program = xelatex
% =====================================================================
% CENTER FOR AI CIVIC KNOWLEDGE — CANONICAL THAI PREPRINT TEMPLATE
% ใช้แม่แบบนี้เป็นฐานสำหรับ paper ภาษาไทยของศูนย์ฯ
% แก้เฉพาะ metadata + เนื้อหาต้นฉบับด้านล่าง
% ห้ามย้ายตำแหน่งโลโก้ / PREPRINT / front matter / หน้า Methodology / header-footer
% =====================================================================

\documentclass{aick-thai-preprint}

% ---------------------------------------------------------------------
% 1) METADATA — EDIT HERE
% ---------------------------------------------------------------------
\AICKTitle{ชื่อบทความภาษาไทย}
\AICKSubtitle{ชื่อรองหรือประโยคที่ระบุขอบเขตของบทความอย่างกระชับ}
\AICKAuthor{ชื่อผู้เขียน (English Name)}
\AICKAuthorShort{นามสกุลผู้เขียน}
\AICKORCID{0000-0000-0000-0000}
\AICKAffiliation{Center for AI Civic Knowledge}
\AICKShortTitle{ชื่อเรื่องย่อสำหรับหัวกระดาษ}
\AICKVersion{v1.0}
\AICKDate{20 กันยายน 2026}

% บทคัดย่อ: แนะนำประมาณ 1,400–1,700 ตัวอักษรไทย
% ให้ผู้อ่านเห็นในครั้งเดียว: ปัญหา → ข้อเสนอ/นิยาม → กลไก → ข้อสรุปหลัก
\AICKAbstract{%
บทคัดย่อควรเริ่มจากปัญหาหรือคำถามที่บทความพยายามตอบ จากนั้นระบุข้อเสนอหลักของบทความให้ชัดเจนโดยไม่ใช้พื้นที่ไปกับการทบทวนวรรณกรรมมากเกินจำเป็น อธิบายกลไกหรือความสัมพันธ์สำคัญที่ทำให้ข้อเสนอนั้นทำงาน และจบด้วยสิ่งที่บทความช่วยให้เราเห็น วัด หรือทดสอบได้แตกต่างจากเดิม

หากบทความมีนิยามแกน ให้ใส่นิยามนั้นในบทคัดย่อโดยตรง แต่หลีกเลี่ยงการอธิบายศัพท์ย่อยทุกคำ รายละเอียดของวรรณกรรม วิธีตรวจ ข้อจำกัด และกรณีตัวอย่างควรย้ายไปอยู่ในเนื้อหาหลัก เพื่อให้หน้าแรกยังคงอ่านได้รวดเร็วและเห็น contribution ของงานในครั้งเดียว
}

\AICKKeywords{คำสำคัญหนึ่ง; คำสำคัญสอง; คำสำคัญสาม; คำสำคัญสี่; คำสำคัญห้า}

\begin{document}

% ---------------------------------------------------------------------
% PAGE 1 — LOCKED ARXIV-LIKE FRONT MATTER
% ---------------------------------------------------------------------
\AICKMakeFrontMatter

% ---------------------------------------------------------------------
% PAGE 2 — REQUIRED METHODOLOGICAL / STATUS NOTE
% เขียนสถานะ paper, จุดกำเนิดของปัญหา, บทบาทมนุษย์–AI,
% วิธีตรวจ/ท้าทายข้อเสนอ, ความรับผิดชอบ และสถานะ PREPRINT
% ---------------------------------------------------------------------
\AICKMethodologyPage{%
บทความนี้เป็นบทความเชิงทฤษฎีและกรอบแนวคิดที่พัฒนาจากปัญหา จุดตั้งต้นคือประสบการณ์หรือข้อสังเกตที่ทำให้เกิดคำถาม ก่อนจะถูกพัฒนาเป็นสมมติฐาน ข้อเสนอ และโครงสร้างที่เปิดให้ตรวจ แก้ หรือปฏิเสธได้ ฉบับปัจจุบันมีสถานะเป็น \textbf{PREPRINT} และยังไม่ผ่านการประเมินโดยผู้ทรงคุณวุฒิ

ในกรณีที่ใช้ปัญญาประดิษฐ์ในกระบวนการพัฒนางาน ให้ระบุบทบาทอย่างตรงไปตรงมา เช่น การช่วยขยายพื้นที่ของการสืบค้น เสนอถ้อยคำ ทางเลือก การเชื่อมโยงข้ามสาขา ข้อคัดค้าน การจัดโครง หรือการตรวจความสอดคล้อง โดยมนุษย์ยังเป็นผู้ถือปัญหา ตัดสินใจเชิงแนวคิด เลือกและตีความหลักฐาน และรับผิดชอบต่อข้อเสนอทั้งหมดของบทความ

ผลจากการสนทนากับ AI ไม่ควรถูกถือเป็นความรู้ที่ได้รับการรับรองโดยอัตโนมัติ แต่เป็นสิ่งที่ควรพิจารณาและต้องผ่านการตรวจด้วยวรรณกรรม หลักฐาน วิธีการอื่น ผู้เชี่ยวชาญ หรือผลจากโลกจริงตามความเหมาะสม ควรแยกให้ชัดว่าอะไรคือประสบการณ์ อะไรคือการตีความ อะไรคือข้อเสนอ และอะไรคือหลักฐานที่ใช้รองรับข้ออ้าง

งานของศูนย์ฯ ยึดหลักวิทยาศาสตร์เปิด การเปิดเผยที่มาของข้อเสนอ การรักษาร่องรอยของการพัฒนาความคิด และการเปิดให้ผู้อื่นตรวจสอบหรือโต้แย้งได้ โดยไม่ใช้สถานะของสถาบัน เครื่องมือ หรือผู้เขียนเป็นสิ่งแทนหลักฐาน
}

% ---------------------------------------------------------------------
% PAGE 3 — MAIN TEXT MUST BEGIN ON A NEW PAGE
% ---------------------------------------------------------------------
\section{บทนำ}
เขียนปัญหาที่เกิดขึ้นในโลกจริงก่อน แล้วระบุคำถามหลักของบทความ หน่วยวิเคราะห์ และ contribution ให้ชัดเจน ผู้อ่านควรเข้าใจภายในช่วงต้นของบทนำว่า “ปัญหาคืออะไร — งานนี้เสนออะไร — และทำไมข้อเสนอนั้นจึงสำคัญ”

\section{กรอบแนวคิดหรือข้อเสนอหลัก}
วางนิยามและกลไกของบทความเองก่อน วรรณกรรมภายนอกควรตามมาในฐานะสิ่งที่ช่วยเปรียบเทียบ ท้าทาย หรือทำให้ข้อเสนอวัดได้ ไม่ควรทำหน้าที่แทนแกนทฤษฎีของบทความหากงานเป็นข้อเสนอเชิงแนวคิดของศูนย์ฯ

\subsection{กลไกที่หนึ่ง}
อธิบายกลไกอย่างเป็นเหตุเป็นผล พร้อมระบุสิ่งที่สังเกตได้หากกลไกนี้เป็นจริง

\subsection{กลไกที่สอง}
แยกความสัมพันธ์ที่มักถูกยุบรวมกัน และระบุเงื่อนไขที่ข้อเสนอนี้ใช้ไม่ได้

\section{ตำแหน่งเมื่อเทียบกับวรรณกรรมที่เกี่ยวข้อง}
ทบทวนวรรณกรรมหลังจากแกนของบทความถูกวางแล้ว โดยใช้ภาษาประเภท “มาบรรจบ”, “ช่วยอธิบาย”, “ต่างตรงที่”, “ทำให้ตรวจได้” มากกว่าการใช้วรรณกรรมเป็นผู้อนุญาตให้ข้อเสนอมีสถานะ

\section{ข้อเสนอเชิงทฤษฎีที่ตรวจสอบได้}
หากเหมาะสม ให้จัดข้อเสนอเป็น P1, P2, P3 หรือเป็นกลุ่มสมมติฐานที่ข้อมูลสามารถสนับสนุน จำกัด หรือหักล้างได้

\section{นัยต่อการประยุกต์ใช้}
อธิบายว่ากรอบนี้เปลี่ยนวิธีออกแบบการเรียนรู้ การทำงาน นโยบาย หรือการตรวจหลักฐานอย่างไร โดยรักษาความต่างระหว่างข้อเสนอเชิงทฤษฎีกับคำแนะนำเชิงปฏิบัติ

\section{ขอบเขต เงื่อนไข และข้อจำกัด}
ระบุให้ชัดว่ากรอบนี้อธิบายอะไรไม่ได้ เงื่อนไขใดทำให้ใช้ไม่ได้ และหลักฐานแบบใดจะทำให้ต้องแก้หรือละทิ้งข้อเสนอ

\section{บทสรุป}
กลับไปตอบคำถามตั้งต้นโดยไม่เพิ่มข้อเสนอใหม่ สรุปกลไกแกน สิ่งที่เปลี่ยนจากความเข้าใจเดิม และคำถามที่ยังเปิดอยู่

% ---------------------------------------------------------------------
% OPTIONAL END MATTER — KEEP ORDER CONSISTENT
% ---------------------------------------------------------------------
\AICKAcknowledgements{%
ระบุเฉพาะบุคคลหรือทรัพยากรที่ควรได้รับการขอบคุณ และไม่ใช้ส่วนนี้แทนการเปิดเผยวิถีวิทยาหรือบทบาท AI ที่ควรอยู่หน้า 2
}

\AICKDeclarations{%
\textbf{สถานะบทความ} PREPRINT; ยังไม่ผ่าน peer review.\par
\textbf{ผลประโยชน์ทับซ้อน} [ระบุหรือเขียนว่าไม่มี].\par
\textbf{ข้อมูลและวัสดุวิจัย} [ระบุแหล่งข้อมูล/โค้ด/คลังเปิด หากมี].
}

\begin{AICKReferences}
\item ผู้แต่ง. (ปี). ชื่อบทความ. \textit{ชื่อวารสาร, เล่ม}(ฉบับ), หน้า. \url{https://doi.org/...}
\item ผู้แต่ง. (ปี). \textit{ชื่อหนังสือ}. สำนักพิมพ์.
\end{AICKReferences}

\end{document}
